\section{Plano de ingeniería: conectar el modelo de razonamiento a un sistema de producción}

\subsection{Arquitectura en capas: Planner / Executor / Verifier}

Combinando lo que se discutió en la entrevista sobre tareas verificables, confiabilidad en cadenas largas y restricciones de costo, un plano de ingeniería práctico es una arquitectura de tres capas:
\begin{itemize}
    \item \textbf{Planner}: responsable de descomponer la tarea, elegir herramientas y controlar el presupuesto
    \item \textbf{Executor}: invoca el modelo y las herramientas para ejecutar las subtareas
    \item \textbf{Verifier}: aplica verificación por reglas, por pruebas o por consistencia a cada producto intermedio
\end{itemize}

\begin{knowledgebox}{Por qué la arquitectura de tres capas es más estable que un ``modelo único que hace todo''}
Un flujo de extremo a extremo con un solo modelo se desarrolla más rápido, pero es difícil controlar los límites de la falla. Al introducir un Verifier, los errores se pueden interceptar en el nivel de la subtarea, en lugar de dejar que se propaguen hasta el resultado final. En tareas de más de 30 pasos, esta diferencia se amplifica de forma exponencial.
\end{knowledgebox}

En este plano, el modelo de razonamiento fuerte debe usarse de manera prioritaria en los pasos de ``alta incertidumbre y alto costo''; los modelos de bajo costo se encargan de los pasos de ``alta frecuencia y bajo riesgo''. Así se mantiene la calidad y a la vez se controla el presupuesto total.

\subsection{Contexto y memoria: evitar que ``entre más se avanza, más se desordena''}

La causa común de que fallen las tareas largas no es que el modelo no sepa hacerlas, sino la contaminación del contexto: la información errónea de etapas tempranas se toma como verdad y se sigue propagando. La discusión en la entrevista sobre el KV Cache y el costo del contexto largo muestra que un ``contexto infinito'' no resuelve por sí solo el problema de la calidad de la memoria.

\begin{table}[H]
\centering
\begin{tabular}{lp{8.7cm}}
\toprule
\textbf{Problema} & \textbf{Medida de ingeniería} \\
\midrule
Expansión del contexto & Usar resúmenes por etapas e instantáneas de estado, en lugar de anexar el historial sin límite \\
Consolidación de memoria errónea & Asignar niveles de confianza a los hechos clave; la información de baja confianza debe verificarse una segunda vez \\
Deriva semántica entre herramientas & Usar una representaci��n intermedia estructurada (JSON schema / estado tipado) \\
Contaminación por reintentos & Cada reintento debe reproducirse desde un estado limpio, para evitar acumular parches sobre una trayectoria errónea \\
\bottomrule
\end{tabular}
\caption{Estrategias de gobernanza del contexto en tareas largas}
\end{table}

\subsection{Estrategia de despliegue: lanzamiento gradual, reversión y división de trabajo entre humano y máquina}

El despliegue del modelo de razonamiento no debe hacerse como una ``sustitución total'', sino como un ``lanzamiento gradual por escenario''. Puede clasificarse por nivel de riesgo:
\begin{itemize}
    \item L1 (riesgo bajo): puede ejecutarse de forma automática, con revisión humana por muestreo
    \item L2 (riesgo medio): el modelo sugiere y el humano confirma
    \item L3 (riesgo alto): el modelo solo da el análisis, sin ejecutar directamente
\end{itemize}

\begin{warningbox}{La forma más peligrosa de desplegar: tratar una tarea L3 como si fuera L1}
Cuando una organización, bajo presión de indicadores clave de desempeño, busca prematuramente una alta tasa de automatización, el error más común es ejecutar de forma automática y directa un proceso de alto riesgo. Cuando se activa este tipo de error, la consecuencia no suele ser una ``caída de precisión'', sino un ``incidente'' de gravedad. La preocupación que muestra la entrevista por los ataques cibernéticos y la vulnerabilidad del sistema eléctrico es, en esencia, un recordatorio de este límite.
\end{warningbox}

Además, la reversión debe ser algo que ``se pueda ensayar'', no algo que solo quede escrito en la documentación. Se recomienda hacer un simulacro de falla una vez al mes: inyectar de forma aleatoria una salida anómala del modelo y verificar si el sistema puede volver a la ruta manual dentro de un tiempo límite.

\subsection{Resumen de la sección}

Llevar el modelo de razonamiento a producción no es solo ingeniería de prompts, sino ingeniería de sistemas. La arquitectura en capas, la gobernanza del contexto, el despliegue por niveles de riesgo y la posibilidad de ensayar la reversión determinan si la capacidad del modelo puede convertirse en productividad estable.
